\section{Related Work}
\label{sec:related}

\paragraph{Hyperspectral Image Classification.}
HSI classification has evolved from pixel-wise spectral methods to deep spatial-spectral approaches. Convolutional networks~\cite{chen2016deep3dcnn}, graph neural networks, and attention-based models have achieved strong performance on standard benchmarks. Recent advances integrate LiDAR elevation data with HSI for multimodal classification, using dual-stream architectures~\cite{xu2024mivit} or cross-attention fusion~\cite{hong2022multimodal}. State-space models (Mamba) have emerged as efficient alternatives to transformers for long-sequence spectral processing~\cite{yao2024spectralgpt}. Our \backbonename{} backbone combines 1D Mamba for spectral analysis with 2D VSSBlocks for spatial processing in a Siamese dual-modal design.

\paragraph{Class-Incremental Learning.}
CIL methods fall into three categories: (i)~\emph{regularization-based} methods like EWC~\cite{kirkpatrick2017ewc} and LwF~\cite{li2017lwf} that constrain parameter updates; (ii)~\emph{replay-based} methods like iCaRL~\cite{rebuffi2017icarl} and LUCIR~\cite{hou2019lucir} that store exemplars; and (iii)~\emph{architecture-based} methods like FOSTER~\cite{wang2022foster} that expand model capacity. For HSI, DCPRN~\cite{deng2024dcprn} uses dual knowledge distillation with prototype replay, CrossACL~\cite{xu2025crossacl} applies analytic continual learning, and PSCEN~\cite{pscen2025} employs prototype similarity calibration. However, all existing HSI-CIL methods operate on single datasets without cross-domain evaluation, and none address multimodal HSI+LiDAR input.

\paragraph{Parameter-Efficient Fine-Tuning for Continual Learning.}
PEFT methods adapt pre-trained models with minimal parameter overhead. Prompt-based approaches (L2P~\cite{wang2022l2p}, DualPrompt~\cite{wang2022dualprompt}, CODA-Prompt~\cite{smith2023codaprompt}) learn task-specific prompts for frozen vision transformers. LoRA-based methods (InfLoRA~\cite{liang2024inflora}, O-LoRA, SD-LoRA~\cite{sdlora2025}) apply low-rank updates with orthogonal or decoupled constraints. EASE~\cite{zhou2024ease} uses expandable adapters with prototype complement. These methods assume a single-modality pre-trained backbone (typically ImageNet ViT) and apply uniform adaptation across all layers. In contrast, \methodname{} leverages the multimodal structure of HSI+LiDAR to make \emph{modality-asymmetric} adaptation decisions---freezing the spectral stream while adapting spatial streams with drift-proportional capacity.

\paragraph{Remote Sensing Continual Learning.}
Cross-domain continual learning for remote sensing has received limited attention. AMoED~\cite{amoed2026} explores domain-specific experts for RS incremental learning but without HSI+LiDAR multimodal input. PMPT applies meta-prompt tuning for multimodal RS few-shot CIL with a fully frozen backbone. Our work differs in three aspects: (i)~we propose a marathon CIL protocol spanning multiple geographic domains with different sensors; (ii)~we selectively freeze only the spectral branch while adapting spatial branches; and (iii)~we provide the first systematic study of cross-domain CIL for multimodal HSI+LiDAR classification.
